% GENERATED FILE. Edit canonical lessons, then run npm run generate:books.
% canonical-source-hash: fnv1a64:24202e95ba5bd418
% canonical-lessons: MW-C184-okhli, MW-C184-chauko, MW-C184-pital, MW-C184-tabo, MW-C184-loho, MW-R184-review-179-184

\chapter{Mortar and Metal}
\label{ch:mw-mortar-and-metal}
\hlchaptermodality{184}

\begin{chapteropening}
\textbf{By the end of this chapter:} \emph{I can name the mortar and the metals pots are made of.}

Complete the last lesson of chapter 184: I can name the mortar and the metals pots are made of.
\end{chapteropening}

\section[\texorpdfstring{\mw{ओखळी} (okhḷī)}{ओखळी (okhḷī)}]{\mw{ओखळी} (okhḷī) — a mortar}

\label{lesson:MW-C184-okhli}



\begin{quote}
Before the new one: say the Marwadi for kohl, then the Marwadi for gram flour.
\end{quote}

\subsection*{You'll want to know: \mw{ओखळी}}
\textbf{\mw{ओखळी}} — \emph{okhḷī} — \textquotedblleft{}a mortar\textquotedblright{}.

\begin{quote}
\mw{ओखळी} \mw{मांय} \mw{मिरच} \mw{है।}
\end{quote}

\begin{quote}
\emph{okhḷī mā̃y mirach hai.}
\end{quote}

\begin{quote}
There are chillies in the mortar.
\end{quote}

A feminine noun.

\begin{cousinweb}
\textbf{\mw{ओखळी}} (\emph{okhḷī}) is inherited from Sanskrit \emph{ulūkhala}, \textquotedblleft{}a wooden mortar\textquotedblright{}; CDIAL lists Hindi \emph{okhlī}, \textquotedblleft{}a small wooden mortar\textquotedblright{} (CDIAL 2360). CDIAL gives no Marwari form in this entry. Its pestle is the \textbf{\mw{मूसळ}} you know.
\end{cousinweb}

\begin{grammarlens}[title={What you've built}]
One more word for this chapter.
\end{grammarlens}

\subsection*{Your turn}
\begin{itemize}
\raggedright
  \item \emph{Say it:} \emph{okhḷī}
  \item \emph{Say it:} \emph{okhḷī}, once more
  \item \emph{Say it:} \emph{okhḷī mā̃y mirach hai.} — There are chillies in the mortar.
  \item \emph{Recall:} say the Marwadi for kohl, then the Marwadi for gram flour, then say \emph{okhḷī} again
\end{itemize}

\begin{tcolorbox}[breakable,colback=teal!4,colframe=teal!35!black,title={Before you move on}]
What does \mw{ओखळी} mean? (\textquotedblleft{}A mortar\textquotedblright{}.) Say it once more.
\end{tcolorbox}
\section[\texorpdfstring{\mw{चौको} (chauko)}{चौको (chauko)}]{\mw{चौको} (chauko) — the cooking place}

\label{lesson:MW-C184-chauko}



\begin{quote}
Before the new one: say the Marwadi for gram flour, then the Marwadi for a mortar.
\end{quote}

\subsection*{You'll want to know: \mw{चौको}}
\textbf{\mw{चौको}} — \emph{chauko} — \textquotedblleft{}the cooking place\textquotedblright{}.

\begin{quote}
\mw{चौको} \mw{साफ} \mw{है।}
\end{quote}

\begin{quote}
\emph{chauko sāph hai.}
\end{quote}

\begin{quote}
The cooking place is clean.
\end{quote}

A masculine noun.

\begin{cousinweb}
\textbf{\mw{चौको}} (\emph{chauko}) is inherited from Sanskrit \emph{catuṣka}, \textquotedblleft{}a set of four\textquotedblright{}, also \textquotedblleft{}a square courtyard\textquotedblright{}; CDIAL lists Hindi \emph{caukā}, a place for cooking and eating (CDIAL 4629). CDIAL gives no Marwari form in this entry.
\end{cousinweb}

\begin{grammarlens}[title={What you've built}]
One more word for this chapter.
\end{grammarlens}

\subsection*{Your turn}
\begin{itemize}
\raggedright
  \item \emph{Say it:} \emph{chauko}
  \item \emph{Say it:} \emph{chauko}, once more
  \item \emph{Say it:} \emph{chauko sāph hai.} — The cooking place is clean.
  \item \emph{Recall:} say the Marwadi for gram flour, then the Marwadi for a mortar, then say \emph{chauko} again
\end{itemize}

\begin{tcolorbox}[breakable,colback=teal!4,colframe=teal!35!black,title={Before you move on}]
What does \mw{चौको} mean? (\textquotedblleft{}The cooking place\textquotedblright{}.) Say it once more.
\end{tcolorbox}
\section[\texorpdfstring{\mw{पीतळ} (pītaḷ)}{पीतळ (pītaḷ)}]{\mw{पीतळ} (pītaḷ) — brass}

\label{lesson:MW-C184-pital}



\begin{quote}
Before the new one: say the Marwadi for a mortar, then the Marwadi for the cooking place.
\end{quote}

\subsection*{You'll want to know: \mw{पीतळ}}
\textbf{\mw{पीतळ}} — \emph{pītaḷ} — \textquotedblleft{}brass\textquotedblright{}.

\begin{quote}
\mw{पीतळ} \mw{पीळो} \mw{है।}
\end{quote}

\begin{quote}
\emph{pītaḷ pīḷo hai.}
\end{quote}

\begin{quote}
Brass is yellow.
\end{quote}

A masculine noun.

\begin{cousinweb}
\textbf{\mw{पीतळ}} (\emph{pītaḷ}) is inherited from Sanskrit \emph{pittala}, \textquotedblleft{}brass\textquotedblright{}; CDIAL lists Gujarati \emph{pitaḷ} and Hindi \emph{pītal} (CDIAL 8184). CDIAL gives no Marwari form in this entry.
\end{cousinweb}

\begin{grammarlens}[title={What you've built}]
One more word for this chapter.
\end{grammarlens}

\subsection*{Your turn}
\begin{itemize}
\raggedright
  \item \emph{Say it:} \emph{pītaḷ}
  \item \emph{Say it:} \emph{pītaḷ}, once more
  \item \emph{Say it:} \emph{pītaḷ pīḷo hai.} — Brass is yellow.
  \item \emph{Recall:} say the Marwadi for a mortar, then the Marwadi for the cooking place, then say \emph{pītaḷ} again
\end{itemize}

\begin{tcolorbox}[breakable,colback=teal!4,colframe=teal!35!black,title={Before you move on}]
What does \mw{पीतळ} mean? (\textquotedblleft{}Brass\textquotedblright{}.) Say it once more.
\end{tcolorbox}
\section[\texorpdfstring{\mw{तांबो} (tā̃bo)}{तांबो (tā̃bo)}]{\mw{तांबो} (tā̃bo) — copper}

\label{lesson:MW-C184-tabo}



\begin{quote}
Before the new one: say the Marwadi for the cooking place, then the Marwadi for brass.
\end{quote}

\subsection*{You'll want to know: \mw{तांबो}}
\textbf{\mw{तांबो}} — \emph{tā̃bo} — \textquotedblleft{}copper\textquotedblright{}.

\begin{quote}
\mw{तांबो} \mw{लाल} \mw{है।}
\end{quote}

\begin{quote}
\emph{tā̃bo lāl hai.}
\end{quote}

\begin{quote}
Copper is red.
\end{quote}

A masculine noun.

\begin{cousinweb}
\textbf{\mw{तांबो}} (\emph{tā̃bo}) is inherited from Sanskrit \emph{tāmra}, \textquotedblleft{}dark red; copper\textquotedblright{}; CDIAL lists Hindi \emph{tā̃bā} (CDIAL 5779). CDIAL gives no Marwari form in this entry.
\end{cousinweb}

\begin{grammarlens}[title={What you've built}]
One more word for this chapter.
\end{grammarlens}

\subsection*{Your turn}
\begin{itemize}
\raggedright
  \item \emph{Say it:} \emph{tā̃bo}
  \item \emph{Say it:} \emph{tā̃bo}, once more
  \item \emph{Say it:} \emph{tā̃bo lāl hai.} — Copper is red.
  \item \emph{Recall:} say the Marwadi for the cooking place, then the Marwadi for brass, then say \emph{tā̃bo} again
\end{itemize}

\begin{tcolorbox}[breakable,colback=teal!4,colframe=teal!35!black,title={Before you move on}]
What does \mw{तांबो} mean? (\textquotedblleft{}Copper\textquotedblright{}.) Say it once more.
\end{tcolorbox}
\section[\texorpdfstring{\mw{लोहो} (loho)}{लोहो (loho)}]{\mw{लोहो} (loho) — iron}

\label{lesson:MW-C184-loho}



\begin{quote}
Before the new one: say the Marwadi for brass, then the Marwadi for copper.
\end{quote}

\subsection*{You'll want to know: \mw{लोहो}}
\textbf{\mw{लोहो}} — \emph{loho} — \textquotedblleft{}iron\textquotedblright{}.

\begin{quote}
\mw{लोहो} \mw{भारी} \mw{है।}
\end{quote}

\begin{quote}
\emph{loho bhārī hai.}
\end{quote}

\begin{quote}
Iron is heavy.
\end{quote}

A masculine noun.

\begin{cousinweb}
\textbf{\mw{लोहो}} (\emph{loho}) is inherited from Sanskrit \emph{loha}, first \textquotedblleft{}copper\textquotedblright{}, later \textquotedblleft{}iron\textquotedblright{}; CDIAL lists Hindi \emph{lohā} (CDIAL 11158). CDIAL gives no Marwari form in this entry.
\end{cousinweb}

\begin{grammarlens}[title={What you've built}]
One more word for this chapter.
\end{grammarlens}

\subsection*{Your turn}
\begin{itemize}
\raggedright
  \item \emph{Say it:} \emph{loho}
  \item \emph{Say it:} \emph{loho}, once more
  \item \emph{Say it:} \emph{loho bhārī hai.} — Iron is heavy.
  \item \emph{Recall:} say the Marwadi for brass, then the Marwadi for copper, then say \emph{loho} again
\end{itemize}

\begin{tcolorbox}[breakable,colback=teal!4,colframe=teal!35!black,title={Before you move on}]
What does \mw{लोहो} mean? (\textquotedblleft{}Iron\textquotedblright{}.) Say it once more.
\end{tcolorbox}
\section[(dialogue)]{Looking back: from Grain and Spice to Mortar and Metal}

\label{lesson:MW-R184-review-179-184}



\begin{quote}
Say the Marwadi for a grinding stone and to pound.
\end{quote}

\subsection*{Your turn}
Say these aloud:

\begin{itemize}
\raggedright
  \item a grinding stone, to pound, to grind, to peel, to winnow
  \item to set down, to put in, to throw, a hanging net, a flat basket
  \item to set curd, to churn, a churning stick, to strain, rice pudding
  \item to shake out, to dissolve, to mix, to temper, to serve food
  \item a ladle, a basket, a shoulder cloth, kohl, gram flour
  \item a mortar, the cooking place, brass, copper, iron
\end{itemize}

\begin{tcolorbox}[breakable,colback=teal!4,colframe=teal!35!black,title={Before you move on}]
A grinding stone? (\textbf{\mw{सिल}}, \emph{sil}.) To pound? (\textbf{\mw{कूटणो}}, \emph{kūṭṇo}.) To grind? (\textbf{\mw{दळणो}}, \emph{daḷṇo}.) To peel? (\textbf{\mw{छीलणो}}, \emph{chhīlṇo}.) To winnow? (\textbf{\mw{फटकणो}}, \emph{phaṭakṇo}.) To set down? (\textbf{\mw{धरणो}}, \emph{dharṇo}.) To put in? (\textbf{\mw{घालणो}}, \emph{ghālṇo}.) To throw? (\textbf{\mw{नाखणो}}, \emph{nākhṇo}.) A hanging net? (\textbf{\mw{छींको}}, \emph{chhī̃ko}.) A flat basket? (\textbf{\mw{छाबड़ी}}, \emph{chhābṛī}.) To set curd? (\textbf{\mw{जमावणो}}, \emph{jamāvṇo}.) To churn? (\textbf{\mw{बिलोवणो}}, \emph{bilovṇo}.) A churning stick? (\textbf{\mw{मथाणी}}, \emph{mathāṇī}.) To strain? (\textbf{\mw{गाळणो}}, \emph{gāḷṇo}.) Rice pudding? (\textbf{\mw{खीर}}, \emph{khīr}.) To shake out? (\textbf{\mw{झटकणो}}, \emph{jhaṭakṇo}.) To dissolve? (\textbf{\mw{घोळणो}}, \emph{ghoḷṇo}.) To mix? (\textbf{\mw{मिलावणो}}, \emph{milāvṇo}.) To temper? (\textbf{\mw{बघारणो}}, \emph{baghārṇo}.) To serve food? (\textbf{\mw{पुरसणो}}, \emph{purasṇo}.) A ladle? (\textbf{\mw{कड़छी}}, \emph{kaṛchhī}.) A basket? (\textbf{\mw{टोकरी}}, \emph{ṭokrī}.) A shoulder cloth? (\textbf{\mw{अंगोछो}}, \emph{aṅgochho}.) Kohl? (\textbf{\mw{काजळ}}, \emph{kājaḷ}.) Gram flour? (\textbf{\mw{बेसण}}, \emph{besaṇ}.) A mortar? (\textbf{\mw{ओखळी}}, \emph{okhḷī}.) The cooking place? (\textbf{\mw{चौको}}, \emph{chauko}.) Brass? (\textbf{\mw{पीतळ}}, \emph{pītaḷ}.) Copper? (\textbf{\mw{तांबो}}, \emph{tā̃bo}.) Iron? (\textbf{\mw{लोहो}}, \emph{loho}.)
\end{tcolorbox}
